%%═══════════════════════════════════════════════════════════════════
%% Lecture 10 — Nuclear Reactions, Cross Sections & Gamma Spectroscopy
%% 77603 — Nuclear Physics
%% Date: 16 June 2026
%%═══════════════════════════════════════════════════════════════════

\section{Lecture 10 — Nuclear Reactions, Cross Sections \& Gamma Spectroscopy}
\label{sec:lec10}
\date{16 June 2026}

\noindent\textbf{\textcolor{doc}{Lecture date: 16 June 2026}}

%%───────────────────────────────────────────────────────────────────
\subsection*{Overview}
%%───────────────────────────────────────────────────────────────────
This lecture moves from static nuclear properties and spontaneous decays to
\emph{nuclear reactions}: controlled collisions in which a projectile is sent
into a target and the outgoing products are measured. Reactions are the main
tool by which we produce unstable nuclei, probe level schemes, measure spins and
parities, and reconstruct astrophysical processes such as stellar burning and
explosive nucleosynthesis.

The first part of the lecture sets up the language of reactions:
$X(a,b)Y$, conservation laws, $Q$-values, threshold energies, laboratory versus
centre-of-mass kinematics, and cross sections. The second part classifies
reaction mechanisms: Coulomb/Rutherford scattering, ordinary nuclear scattering,
compound-nucleus reactions, direct reactions, and resonant reactions described by
the Breit--Wigner formula. The final part is a detailed worked example from
Exercise~3: reconstructing the $\gamma$ spectrum and level feeding in
$\nucleus{Na}{24}\to\nucleus{Mg}{24}$ after $\beta^-$ decay. This example ties
together mass excesses, $\beta$ decay, $\gamma$ spectroscopy, detector response,
branching ratios, coincidence peaks, and spin-parity selection rules.

%%═══════════════════════════════════════════════════════════════════
\subsection{What Is a Nuclear Reaction?}
\label{subsec:lec10:reaction_notation}
%%═══════════════════════════════════════════════════════════════════

\subsubsection*{Physical Motivation}
Until now, most of the course asked what a nucleus \emph{is}: its mass, radius,
binding energy, spin, parity, shell structure, and decay modes. A reaction asks a
more dynamical question: what happens when we deliberately bring two nuclear
systems together? This matters in at least three settings:
\begin{itemize}[nosep]
  \item \textbf{astrophysics}: stars, novae, supernovae, and neutron-rich
        explosions are powered or shaped by nuclear reactions;
  \item \textbf{experiments}: scattering and transfer reactions reveal nuclear
        radii, level energies, spins, parities, and spectroscopic information;
  \item \textbf{applications}: medical isotopes, activation analysis, reactor
        physics, and radiometric methods all depend on controlled reactions.
\end{itemize}
The energy scale in this discussion is nuclear rather than particle-physics
scale: typically a few $\text{MeV}$ to tens of $\text{MeV}$, comparable to
nuclear binding energies. At substantially higher energies, additional degrees
of freedom enter and the simple low-energy reaction language must be extended.

\dfn[dfn:lec10:reaction_notation]{Reaction Notation}{%
A general two-body entrance and exit channel is written
\begin{equation}
  \label{eq:lec10:reaction_long}
  a + X \longrightarrow b + Y,
\end{equation}
or, in the experimental shorthand,
\begin{equation}
  \label{eq:lec10:reaction_shorthand}
  \boxed{X(a,b)Y.}
\end{equation}
Here $X$ is the target nucleus, $a$ is the projectile, $b$ is the measured
outgoing light particle, and $Y$ is the residual nucleus.}
\textit{Significance:} the shorthand $X(a,b)Y$ tells the experimental story:
we started with target $X$, fired projectile $a$, detected particle $b$, and
therefore inferred residual nucleus $Y$.

\paragraph{Examples.}
Elastic scattering is $X(a,a)X$: the same particles appear before and after.
Inelastic scattering is $X(a,a')X^\ast$: the projectile remains the same species
but loses kinetic energy to excite the target. A transfer reaction such as
$\nucleus{Li}{7}(p,n)\nucleus{Be}{7}$ converts the target by transferring charge
and nucleons. A radiative capture reaction such as
\begin{equation}
  \label{eq:lec10:radiative_capture_example}
  \nucleus{Au}{197}(n,\gamma)\nucleus{Au}{198}
\end{equation}
absorbs the projectile and removes the excess energy by photon emission.

\nt{The notation is conventional, not ontological. The same centre-of-mass
reaction can sometimes be studied in inverse kinematics: instead of firing a
light projectile at a heavy target, one accelerates the heavy radioactive nucleus
and sends it into a light target. This is essential when the nucleus of interest
is too short-lived to be made into a stationary target.}

%%═══════════════════════════════════════════════════════════════════
\subsection{Conservation Laws and Q-Values}
\label{subsec:lec10:conservation_q}
%%═══════════════════════════════════════════════════════════════════

\subsubsection*{Conceptual Background}
A nuclear reaction is fast, with a characteristic interaction time of order
$10^{-22}$--$10^{-16}\,\text{s}$ depending on the mechanism. On this time scale
the strong and electromagnetic interactions dominate; weak $\beta$ conversion is
not part of the reaction itself. Therefore the reaction must obey the ordinary
conservation laws, but with an especially useful simplification: proton number
and neutron number are conserved separately in the reaction stage.

\thm[th:lec10:conservation]{Conservation Laws in Low-Energy Nuclear Reactions}{%
For $a+X\to b+Y$ one must conserve
\begin{align}
  \label{eq:lec10:conserve_energy_momentum}
  E_\text{tot}^{(i)} &= E_\text{tot}^{(f)}, &
  \bvec p_\text{tot}^{(i)} &= \bvec p_\text{tot}^{(f)},\\
  \label{eq:lec10:conserve_quantum_numbers}
  Z_i &= Z_f, &
  N_i &= N_f,\\
  \label{eq:lec10:conserve_spin_parity}
  J_i^\pi &\to J_f^\pi \quad\text{subject to angular-momentum and parity
  coupling.}
\end{align}}
\textit{Significance:} separate conservation of $Z$ and $N$ makes reaction
bookkeeping much sharper than merely conserving $A=Z+N$; $\beta$ decay may occur
later, but it is not part of the prompt nuclear reaction.

\subsubsection*{Mathematical Setup --- the $Q$-Value}
The $Q$-value is the rest-mass energy released by the reaction:
\begin{equation}
  \label{eq:lec10:q_value}
  \boxed{
  Q = \bigl(M_a + M_X - M_b - M_Y\bigr)c^2.}
\end{equation}
Equivalently, by total-energy conservation,
\begin{equation}
  \label{eq:lec10:q_value_kinetic}
  \boxed{
  Q = T_b + T_Y - T_a - T_X,}
\end{equation}
where $T_i$ denotes kinetic energy. If the target is at rest in the laboratory,
$T_X=0$.
\textit{Significance:} $Q>0$ means the reaction is exothermic; $Q<0$ means the
incoming beam must supply at least enough kinetic energy to make up the rest-mass
deficit.

\subsubsection*{Threshold Energy}
For an endothermic reaction, the threshold is \emph{not} simply
$T_a=-Q$. Some kinetic energy must remain as centre-of-mass motion because total
momentum must be conserved. The clean way to see this is to use the
centre-of-mass energy. For a projectile $a$ incident on a stationary target $X$,
the available centre-of-mass kinetic energy is
\begin{equation}
  \label{eq:lec10:e_cm_lab}
  E_{\text{cm}} =
  T_a\,\frac{M_X}{M_a+M_X}.
\end{equation}
At threshold the products have no relative kinetic energy in the centre-of-mass
frame, so all available centre-of-mass energy is used to pay $-Q$:
\begin{equation}
  \label{eq:lec10:threshold_derivation}
  T_{a,\text{thr}}\,\frac{M_X}{M_a+M_X} = -Q.
\end{equation}
Therefore
\begin{equation}
  \label{eq:lec10:threshold}
  \boxed{
  T_{a,\text{thr}} =
  -Q\,\frac{M_a+M_X}{M_X}
  \qquad(Q<0,\; X\text{ at rest}).}
\end{equation}
\textit{Significance:} the projectile must provide more than $|Q|$ because the
final system must still carry the original centre-of-mass motion.

\subsubsection*{Lab and Centre-of-Mass Pictures}
The same reaction has a simpler geometry in the centre-of-mass frame: total
momentum is zero, so the outgoing products are back-to-back. To compare with the
detector, one then adds the centre-of-mass velocity vector to each outgoing
velocity. This vector addition has an important experimental consequence. If the
outgoing particle's centre-of-mass speed is smaller than the centre-of-mass
velocity, the particle always moves forward in the lab, and the same laboratory
angle can correspond to two centre-of-mass emission directions. If the outgoing
speed is larger, backward laboratory emission becomes possible and the mapping
changes.

\nt{This is why reaction kinematics plots can have two branches at the same lab
angle. The ambiguity is not detector noise; it is a real consequence of adding
the centre-of-mass velocity to an isotropic two-body exit channel.}

%%═══════════════════════════════════════════════════════════════════
\subsection{Cross Sections: Probability Written as an Area}
\label{subsec:lec10:cross_sections}
%%═══════════════════════════════════════════════════════════════════

\subsubsection*{Physical Motivation}
The cross section answers the practical question: if I fire a projectile through
a target, what is the probability that the reaction I care about occurs? It has
units of area because of a simple geometric analogy. If a blind shooter fires
uniformly through a frame containing a drawn object, the hit probability is the
object's area divided by the frame area. In a target foil, the ``objects'' are
effective interaction areas attached to each nucleus.

\dfn[dfn:lec10:total_cross_section]{Total Cross Section}{%
For a thin target,
\begin{equation}
  \label{eq:lec10:probability_sigma}
  \boxed{
  \frac{N_\text{react}}{N_\text{in}} = \sigma\,t,}
\end{equation}
where $t$ is the number of target nuclei per unit area:
\begin{equation}
  \label{eq:lec10:areal_density}
  t = \frac{\rho\,\ell\,N_A}{A}.
\end{equation}
Here $\rho$ is mass density, $\ell$ is target thickness, $N_A$ is Avogadro's
number, and $A$ is the molar mass in $\text{g/mol}$.}
\textit{Significance:} $\sigma t$ is dimensionless and is the reaction
probability in the thin-target limit $\sigma t\ll1$.

For a thick target, one must account for removal of particles as they pass
through successive layers. The infinitesimal form
$dI=-n\sigma I\dd{x}$ gives
\begin{equation}
  \label{eq:lec10:attenuation}
  \boxed{I(x)=I_0 e^{-n\sigma x}.}
\end{equation}
\textit{Significance:} the linear probability formula is the first-order limit
of exponential attenuation.

\subsubsection*{Why It Is Not Just a Geometric Size}
If nuclei were hard spheres and the projectile had no wave nature, one might
guess $\sigma=\pi R^2$. Quantum mechanically this is only a rough scale. The
projectile has reduced de Broglie wavelength
\begin{equation}
  \label{eq:lec10:reduced_wavelength}
  \bar\lambda = \frac{\hbar}{p},
\end{equation}
so even a hard-sphere estimate becomes energy dependent:
\begin{equation}
  \label{eq:lec10:hard_sphere}
  \boxed{\sigma_{\text{hard sphere}}\sim \pi\qty(R+\bar\lambda)^2.}
\end{equation}
\textit{Significance:} slow particles have larger wavelengths and can have
larger effective cross sections. More importantly, $\sigma$ also depends on the
specific reaction channel and on selection rules, resonances, and nuclear
structure.

\paragraph{Units.}
The standard nuclear unit is the barn:
\begin{equation}
  \label{eq:lec10:barn}
  \boxed{1\,\text{b}=10^{-28}\,\text{m}^2=10^{-24}\,\text{cm}^2
  =100\,\text{fm}^2.}
\end{equation}
Typical nuclear cross sections range from microbarns for rare processes to barns
or more for favourable channels.

\subsubsection*{Differential Cross Sections}
Experiments often do not ask only whether a reaction occurred. They ask where
the outgoing particle went and with what energy. The corresponding observables
are
\begin{align}
  \label{eq:lec10:diff_cross_sections}
  \frac{d\sigma}{d\Omega}(E_a,\theta)
    &\quad\text{angular distribution},\\
  \frac{d\sigma}{dE_b}(E_a,E_b)
    &\quad\text{outgoing-energy distribution},\\
  \frac{d^2\sigma}{dE_b\dd{\Omega}}(E_a,E_b,\theta)
    &\quad\text{double-differential distribution}.
\end{align}
The total cross section is obtained by integrating over the unobserved variables:
\begin{equation}
  \label{eq:lec10:integral_cross_section}
  \boxed{\sigma(E_a)=\int \frac{d\sigma}{d\Omega}(E_a,\theta)\dd{\Omega}.}
\end{equation}
\textit{Significance:} angular and energy distributions are where the nuclear
structure information sits; the total cross section alone is often too blunt.

%%═══════════════════════════════════════════════════════════════════
\subsection{Reaction Mechanisms}
\label{subsec:lec10:mechanisms}
%%═══════════════════════════════════════════════════════════════════

\subsubsection*{Physical Motivation}
Two reactions with the same entrance and exit particles need not proceed through
the same microscopic story. The mechanism is determined by energy, wavelength,
time scale, and whether the projectile interacts with the whole nucleus or with
individual nucleons.

\subsubsection*{Coulomb and Nuclear Scattering}
At low energies and large distances, charged projectiles primarily feel the
Coulomb field. Then scattering is Rutherford-like: the projectile sees only the
long-range electric force. Once the energy is high enough that the projectile
reaches nuclear distances, the short-range nuclear force contributes and the
wave nature of the projectile produces diffraction/interference patterns.

\thm[th:lec10:rutherford]{Rutherford Scattering Regime}{%
For a point Coulomb field, the differential cross section has the characteristic
angular dependence
\begin{equation}
  \label{eq:lec10:rutherford}
  \boxed{
  \frac{d\sigma}{d\Omega}
  \propto \frac{1}{\sin^4(\theta/2)}.}
\end{equation}}
\textit{Significance:} deviations from the Rutherford form signal finite nuclear
size, electron screening, relativistic/spin effects, or the onset of the nuclear
force.

\subsubsection*{Compound-Nucleus Reactions}
In a compound reaction, the projectile is fully absorbed and its energy is shared
among many nucleons:
\begin{equation}
  \label{eq:lec10:compound_sequence}
  a+X \longrightarrow C^\ast \longrightarrow b+Y.
\end{equation}
The lifetime is long on nuclear collision scales, typically
$10^{-18}$--$10^{-16}\,\text{s}$. That is enough time for the compound nucleus to
lose memory of how it was formed.

\dfn[dfn:lec10:compound]{Compound-Nucleus Reaction}{%
A reaction mechanism in which the entrance channel forms an equilibrated excited
nucleus $C^\ast$, which later decays by emitting particles or photons.}
\textit{Significance:} the exit probabilities depend mainly on the properties of
$C^\ast$ and the energetically open decay channels, not on the detailed entrance
history.

This mechanism is also called a fusion-evaporation process when the two incoming
nuclei fuse and the hot compound nucleus cools by evaporating nucleons or light
clusters. It is useful for isotope production and for studying highly excited
nuclei.

\subsubsection*{Direct Reactions}
At higher energies the projectile wavelength is shorter and the interaction time
is much shorter, of order $10^{-22}\,\text{s}$. The projectile no longer sees a
smooth whole nucleus; it samples individual nucleons or clusters. It may knock
out a nucleon, transfer a nucleon, or excite a specific configuration.

\dfn[dfn:lec10:direct]{Direct Reaction}{%
A fast reaction in which the projectile interacts with one or a few degrees of
freedom in the target, without forming an equilibrated compound nucleus.}
\textit{Significance:} direct reactions are sensitive to single-particle or
cluster structure, so they are powerful spectroscopic probes of shell-model
orbitals and nuclear wavefunctions.

\subsubsection*{Resonant Reactions}
A resonant reaction occurs when the entrance-channel energy matches a real
excited state of the combined system. Then the cross section can become very
large over a narrow energy interval. If the system returns to the entrance
channel, we observe resonant scattering; if it decays by photon emission, we
observe radiative capture:
\begin{equation}
  \label{eq:lec10:resonance_capture}
  a+X \longrightarrow C^\ast(E_r) \longrightarrow C+\gamma.
\end{equation}

\dfn[dfn:lec10:partial_width]{Partial Widths}{%
If a resonant state can decay through several channels $i$, each channel has a
partial width
\begin{equation}
  \label{eq:lec10:partial_width}
  \Gamma_i = \frac{\hbar}{\tau_i},
\end{equation}
and the total width is
\begin{equation}
  \label{eq:lec10:total_width}
  \boxed{\Gamma=\sum_i \Gamma_i.}
\end{equation}}
\textit{Significance:} $\Gamma_i/\Gamma$ is the branching fraction for channel
$i$, while $\Gamma$ fixes the total lifetime and the energy width of the
resonance.

\thm[th:lec10:breit_wigner]{Breit--Wigner Resonance Formula}{%
For entrance channel $a+X$ and exit channel $b+Y$, a single isolated resonance
contributes approximately
\begin{equation}
  \label{eq:lec10:breit_wigner}
  \boxed{
  \sigma_{a\to b}(E)
  =
  \frac{\pi}{k^2}\,
  \frac{2J+1}{(2I_a+1)(2I_X+1)}\,
  \frac{\Gamma_a\Gamma_b}
       {(E-E_r)^2+\Gamma^2/4}.}
\end{equation}}
Here $k$ is the entrance-channel wave number, $J$ is the resonance spin,
$I_a,I_X$ are the projectile and target spins, $E_r$ is the resonance energy,
and $\Gamma_a,\Gamma_b$ are the entrance and exit partial widths.
\textit{Significance:} the Lorentzian denominator says \emph{where} the
resonance is and how broad it is; the numerator says that both forming the
resonance and decaying into the observed channel must be probable.

%%═══════════════════════════════════════════════════════════════════
\subsection{Worked Example: Na-24 to Mg-24 Gamma Spectroscopy}
\label{subsec:lec10:na24_exercise}
%%═══════════════════════════════════════════════════════════════════

\subsubsection*{Problem Statement}
Exercise~3 studies the decay
\begin{equation}
  \label{eq:lec10:na24_decay}
  \nucleus{Na}{24}\longrightarrow \nucleus{Mg}{24}+e^-+\bar\nu_e,
\end{equation}
followed by $\gamma$ de-excitation of excited $\nucleus{Mg}{24}$ levels. The
given $\nucleus{Mg}{24}$ levels are
\begin{center}
\begin{tabular}{ccl}
  Energy (keV) & $J^\pi$ & label \\
  \hline
  $0$       & $0^+$ & ground state \\
  $1368.67$ & $2^+$ & first excited state \\
  $4122.85$ & $4^+$ & strongly fed level \\
  $4238.35$ & $2^+$ & weakly fed/possibly un-fed level \\
  $5235.16$ & $3^+$ & weakly fed level
\end{tabular}
\end{center}
The observed $\gamma$ lines are
\begin{center}
\begin{tabular}{ccc}
  $E_\gamma$ (keV) & $I_\gamma$ (\%) & interpretation \\
  \hline
  $996.82$   & $0.0014$ & $5235.16\to4238.35$ \\
  $1368.633$ & $100$    & $1368.67\to0$ \\
  $2754.028$ & $99.944$ & $4122.85\to1368.67$ \\
  $2869.5$   & $0.0003$ & $4238.35\to1368.67$ \\
  $3866.19$  & $0.052$  & $5235.16\to1368.67$ \\
  $4237.96$  & $0.0011$ & $4238.35\to0$
\end{tabular}
\end{center}
The intensities are normalized to the $1368.633\,\text{keV}$ line.

\subsubsection*{Step 1 --- Compute the $\beta^-$ $Q$-Value}
The exercise gives atomic mass excesses
\begin{equation}
  \label{eq:lec10:mass_excess_values}
  \Delta(\nucleus{Na}{24})=-8417.901(17)\,\text{keV},\qquad
  \Delta(\nucleus{Mg}{24})=-13933.578(13)\,\text{keV}.
\end{equation}
For $\beta^-$ decay using atomic masses, the electron masses cancel:
\begin{equation}
  \label{eq:lec10:beta_minus_q_atomic}
  Q_{\beta^-}
  = \bigl[M_\text{atom}(A,Z)-M_\text{atom}(A,Z+1)\bigr]c^2.
\end{equation}
Since both atoms have the same $A=24$, the $24u$ parts cancel and the mass
excesses alone determine $Q$:
\begin{align}
  \label{eq:lec10:na24_q_calc}
  Q_{\beta^-}
  &= \Delta(\nucleus{Na}{24})-\Delta(\nucleus{Mg}{24}) \nonumber\\
  &= -8417.901 - (-13933.578)\,\text{keV} \nonumber\\
  &= 5515.677\,\text{keV}.
\end{align}
The uncertainty is
\begin{equation}
  \label{eq:lec10:na24_q_uncertainty}
  \delta Q = \sqrt{(0.017)^2+(0.013)^2}\,\text{keV}
  = 0.021\,\text{keV}.
\end{equation}
Therefore
\begin{equation}
  \label{eq:lec10:na24_q_result}
  \boxed{Q_{\beta^-}=5515.677(21)\,\text{keV}.}
\end{equation}
\textit{Significance:} the $Q$-value is large enough to populate all listed
excited states of $\nucleus{Mg}{24}$, but it does not say by itself which levels
are actually populated; that requires spectroscopy and selection rules.

\subsubsection*{Step 2 --- Identify the Spectrum}
Although the parent decay is $\beta^-$, the displayed spectrum is not a
$\beta$ spectrum. A $\beta$ spectrum is continuous because the electron and
antineutrino share energy. The measured spectrum contains sharp peaks plus a
large continuum, the signature of $\gamma$ detection in matter.

\clm[clm:lec10:gamma_spectrum]{The measured spectrum is a gamma spectrum}{%
The decisive clues are:
\begin{itemize}[nosep]
  \item sharp full-energy peaks at transition energies;
  \item Compton continua and Compton edges;
  \item pair-production escape peaks and the $511\,\text{keV}$ annihilation line.
\end{itemize}}
\textit{Significance:} the experiment detects photons emitted by excited
$\nucleus{Mg}{24}$ after the $\beta$ decay, not the $\beta$ electron energy
distribution itself.

\subsubsection*{Step 3 --- Detector Features: Escape Peaks, Annihilation, Compton Edge}
For a photon with energy $E_\gamma>2m_ec^2=1022\,\text{keV}$, pair production can
occur in or near the detector:
\begin{equation}
  \label{eq:lec10:pair_production}
  \gamma \to e^-+e^+.
\end{equation}
The positron annihilates,
\begin{equation}
  \label{eq:lec10:annihilation}
  e^+ + e^- \to 2\gamma_{511},
\end{equation}
producing two $511\,\text{keV}$ photons. If one or both annihilation photons
escape the detector, the deposited energy is
\begin{equation}
  \label{eq:lec10:escape_peaks}
  \boxed{
  E_\text{single escape}=E_\gamma-511\,\text{keV},\qquad
  E_\text{double escape}=E_\gamma-1022\,\text{keV}.}
\end{equation}
\textit{Significance:} not every peak in a $\gamma$ spectrum is a nuclear
transition; some are detector-response peaks.

For the $1368.633\,\text{keV}$ line,
\begin{align}
  \label{eq:lec10:escape_1368}
  E_\text{single} &\simeq 1368.633-511 = 857.6\,\text{keV},\\
  E_\text{double} &\simeq 1368.633-1022 = 346.6\,\text{keV}.
\end{align}
Thus the unexplained $346\,\text{keV}$ feature is naturally the double-escape
peak associated with the $1368.6\,\text{keV}$ photon. For the
$2754.0\,\text{keV}$ line,
\begin{equation}
  \label{eq:lec10:single_escape_2754}
  2754.028-511 \simeq 2243\,\text{keV},
\end{equation}
so the $2243\,\text{keV}$ peak is the single-escape peak of the
$2754.028\,\text{keV}$ transition. The $511\,\text{keV}$ peak itself is the
annihilation line.

The broad background below full-energy peaks comes mainly from Compton
scattering. The scattered photon energy is
\begin{equation}
  \label{eq:lec10:compton_photon}
  E_{\gamma'}=
  \frac{E_\gamma}{1+\dfrac{E_\gamma}{m_ec^2}(1-\cos\theta)},
\end{equation}
so the electron receives
\begin{equation}
  \label{eq:lec10:compton_electron_energy}
  T_e=E_\gamma-E_{\gamma'}.
\end{equation}
The maximum deposited Compton energy occurs for backscattering,
$\theta=\pi$:
\begin{equation}
  \label{eq:lec10:compton_edge}
  \boxed{
  T_{e,\max}
  =
  E_\gamma-
  \frac{E_\gamma}{1+2E_\gamma/(m_ec^2)}.}
\end{equation}
\textit{Significance:} the red-arrow feature in the exercise is a Compton edge:
the sharp upper end of the continuum from photons that scatter once and escape.

\subsubsection*{Step 4 --- Coincidence Sum Peak}
The spectrum also shows a peak near $4122\,\text{keV}$ that is not a listed
nuclear transition. It is naturally explained as a \emph{sum peak}:
\begin{equation}
  \label{eq:lec10:sum_peak}
  \boxed{
  1368.633\,\text{keV}+2754.028\,\text{keV}
  = 4122.661\,\text{keV}.}
\end{equation}
If the $4122.85\,\text{keV}$ level decays by the cascade
\[
  4122.85 \to 1368.67 \to 0,
\]
the two photons can arrive in the detector within the shaping time of the
electronics. The detector then records one event whose energy is the sum of the
two photon energies.

\textit{Significance:} a sum peak is evidence that the two photons are emitted in
prompt cascade from the same decay chain. It is not an extra level by itself.

\subsubsection*{Step 5 --- Why Intensities Do Not Behave Like Peak Heights}
The table's intensities are nuclear line intensities, normalized to the
$1368.633\,\text{keV}$ line. They need not match the raw peak heights in the
displayed spectrum. The observed full-energy peak area is schematically
\begin{equation}
  \label{eq:lec10:observed_peak_area}
  A_\text{peak}(E_\gamma)
  =
  N_\text{decay}\,
  I_\gamma\,
  \varepsilon_\text{full}(E_\gamma),
\end{equation}
where $\varepsilon_\text{full}(E_\gamma)$ is the full-energy peak efficiency.
This efficiency depends strongly on photon energy and on competing processes:
photoelectric absorption, Compton escape, and pair-production escape.

\textit{Significance:} the $1368.6\,\text{keV}$ and $2754.0\,\text{keV}$ lines
can have nearly equal nuclear intensities while their plotted photopeak heights
differ by orders of magnitude, because the detector is much less efficient for
recording the full energy of the higher-energy photon.

The sum of listed line intensities can exceed $100\%$ because one parent decay
can emit more than one photon in a cascade. Intensities count photons, not
exclusive decay branches.

\subsubsection*{Step 6 --- Reconstruct Which Levels Were Populated}
The central spectroscopy logic is a flow balance. For each level,
\begin{equation}
  \label{eq:lec10:level_balance}
  \boxed{
  \text{direct }\beta\text{ feeding}
  =
  \sum I_\gamma(\text{out of level})
  -
  \sum I_\gamma(\text{into level}),}
\end{equation}
after correcting all line intensities for detector efficiency and internal
conversion when needed. In this exercise the table already gives the intensities
to use, so the qualitative balance can be done directly.
\textit{Significance:} level feeding is not read from a single line; it is a
conservation-of-population problem on the level diagram.

\paragraph{The $4122.85\,\text{keV}$ level.}
The $4122.85\,\text{keV}$ level decays by the strong
$2754.028\,\text{keV}$ transition to the $1368.67\,\text{keV}$ level:
\begin{equation}
  \label{eq:lec10:feed_4122}
  B_\beta(4122.85)\simeq I_\gamma(2754.028)=99.944\%.
\end{equation}
There is no listed strong feeding into this level from above. Therefore almost
all $\nucleus{Na}{24}$ $\beta$ decays populate this $4^+$ level directly.

\paragraph{The $5235.16\,\text{keV}$ level.}
This level decays through two listed branches:
\[
  5235.16\to4238.35 \quad(996.82\,\text{keV}),\qquad
  5235.16\to1368.67 \quad(3866.19\,\text{keV}).
\]
Thus
\begin{equation}
  \label{eq:lec10:feed_5235}
  B_\beta(5235.16)
  \simeq 0.0014\%+0.052\%
  = 0.0534\%.
\end{equation}
\textit{Significance:} this level is populated, but only weakly; the
$4122.85\,\text{keV}$ level remains the dominant branch by far.

\paragraph{The $4238.35\,\text{keV}$ level.}
Outgoing intensity from this level is
\begin{equation}
  \label{eq:lec10:out_4238}
  I_\text{out}(4238.35)
  =
  I_\gamma(2869.5)+I_\gamma(4237.96)
  =
  0.0003\%+0.0011\%
  =
  0.0014\%.
\end{equation}
Incoming intensity from the $5235.16\,\text{keV}$ level is
\begin{equation}
  \label{eq:lec10:in_4238}
  I_\text{in}(4238.35)=I_\gamma(996.82)=0.0014\%.
\end{equation}
Therefore
\begin{equation}
  \label{eq:lec10:feed_4238}
  \boxed{
  B_\beta(4238.35)\simeq 0.0014\%-0.0014\%=0.}
\end{equation}
\textit{Significance:} within the quoted uncertainties, there is no need to
postulate direct $\beta$ feeding of the $4238.35\,\text{keV}$ level. A tiny
feeding cannot be excluded, but the data are consistent with zero.

\paragraph{The $1368.67\,\text{keV}$ level.}
This level emits the normalized $1368.633\,\text{keV}$ photon with intensity
$100\%$. It also receives feeding from above:
\begin{equation}
  \label{eq:lec10:in_1368}
  I_\text{in}(1368.67)
  =
  I_\gamma(2754.028)+I_\gamma(2869.5)+I_\gamma(3866.19).
\end{equation}
Numerically,
\begin{equation}
  \label{eq:lec10:feed_1368}
  B_\beta(1368.67)
  \simeq
  100 - 99.944 - 0.0003 - 0.052
  =
  0.0037\%.
\end{equation}
Given the uncertainties, this is a very small value and is consistent with
essentially no direct feeding at the level of this rough reconstruction.

\subsubsection*{Step 7 --- Spin and Parity of the $\nucleus{Na}{24}$ Ground State}
The strongest observed $\beta$ feeding goes to the $4122.85\,\text{keV}$ state,
whose spin and parity are $4^+$. A strong $\beta$ branch usually indicates an
allowed or nearly allowed transition. Allowed $\beta$ transitions do not change
parity and have
\begin{equation}
  \label{eq:lec10:allowed_beta_selection}
  \Delta J=0 \quad\text{(Fermi)},\qquad
  \Delta J=0,\pm1 \quad\text{(Gamow--Teller, not }0\to0\text{)}.
\end{equation}
The simplest inference is therefore
\begin{equation}
  \label{eq:lec10:na24_spin_parity}
  \boxed{J^\pi\bigl(\nucleus{Na}{24}_\text{g.s.}\bigr)\approx4^+.}
\end{equation}
\textit{Significance:} the dominant feeding acts like a spectroscopic pointer:
the parent ground state is most naturally matched to the strongly fed $4^+$
state of $\nucleus{Mg}{24}$.

This inference agrees with the literature value, $J^\pi(\nucleus{Na}{24})=4^+$.

\subsubsection*{Can We See Feeding to the $\nucleus{Mg}{24}$ Ground State?}
Pure feeding directly into the $\nucleus{Mg}{24}$ ground state would emit no
subsequent $\gamma$ ray, so it cannot be measured from the $\gamma$ intensity
balance alone. However, selection rules provide a strong physical argument. If
the parent is $4^+$ and the daughter ground state is $0^+$, then
\begin{equation}
  \label{eq:lec10:ground_state_delta_j}
  \Delta J=4,\qquad \Delta\pi=+.
\end{equation}
The leptons must carry a large angular momentum, so the transition is highly
forbidden and strongly suppressed compared with the observed $4^+\to4^+$
feeding.

\nt{Exam strategy: when reconstructing a decay scheme, separate what the
spectrum directly proves from what selection rules make plausible. Gamma
intensities cannot prove a silent ground-state branch is absent, but spin-parity
arguments can show it should be extremely small.}

\subsubsection*{Key Insight of the Exercise}
The exercise is not mainly about recognizing individual peaks. The deeper method
is:
\begin{enumerate}[label=(\roman*),nosep]
  \item use mass excesses to establish the available $\beta$ energy;
  \item identify which spectral features are real nuclear transitions and which
        are detector response;
  \item correct the intuition about peak heights by remembering detector
        efficiency;
  \item balance incoming and outgoing $\gamma$ intensities level by level;
  \item use the dominant feeding and selection rules to infer spin and parity.
\end{enumerate}
Knowing this chain is equivalent to knowing a large fraction of the course at
the level required for the final exam.
